\begin{table}[H]
\centering
\caption{Mayores y menores crecimientos del PIB, 61 agrupaciones CIIU, 2010--2025}
\label{tab:valor_agregado_61_agrupaciones}
\small
\begin{tabular}{llp{0.46\textwidth}rr}
\toprule
Grupo & Código & Actividad económica & 2025 & Crec. \\
\midrule
Mayores & 104-108 & Artes, entretenimiento y otros servicios & 40,6 & 9,09\% \\
Mayores & 075 & Transporte aéreo & 7,7 & 6,75\% \\
Mayores & 085-088 & Financieras y seguros & 53,1 & 5,95\% \\
Mayores & 102,103 & Salud y servicios sociales & 49,8 & 5,11\% \\
Mayores & 098,099 & Administración pública y defensa & 62,7 & 4,75\% \\
Mayores & 101 & Educación de no mercado & 30,8 & 4,67\% \\
Mayores & 076 & Almacenamiento y apoyo al transporte & 9,5 & 4,50\% \\
Mayores & 016 & Pesca y acuicultura & 2,6 & 4,44\% \\
Mayores & 070 & Comercio & 82,3 & 3,85\% \\
Mayores & 052 & Equipo eléctrico y electrónico & 2,9 & 3,72\% \\
Menores & 019 & Minerales metalíferos & 3,0 & -2,32\% \\
Menores & 039 & Cuero, calzado y artículos de viaje & 1,2 & -2,28\% \\
Menores & 017 & Carbón & 5,7 & -1,98\% \\
Menores & 022 & Apoyo a otras actividades mineras & 0,1 & -0,75\% \\
Menores & 073 & Transporte acuático & 0,2 & -0,73\% \\
Menores & 067 & Edificaciones & 20,5 & -0,49\% \\
Menores & 040 & Madera & 1,0 & -0,26\% \\
Menores & 018,021 & Petróleo, gas y apoyo conexo & 24,6 & -0,22\% \\
Menores & 053,057 & Maquinaria e instalación & 4,0 & -0,16\% \\
Menores & 020 & Otras minas y canteras & 1,5 & 0,01\% \\
\bottomrule
\end{tabular}
\caption*{\footnotesize Nota: valores de 2025 en billones de pesos constantes de 2015. La tabla muestra los diez mayores y los diez menores crecimientos anualizados entre las 61 agrupaciones de actividad económica del DANE. En sentido estricto, el numerador corresponde al valor agregado bruto de cada actividad. La tabla completa se deja como archivo de resultados del proyecto. Fuente: cálculos propios con DANE, anexo de producción a precios constantes.}
\end{table}

\textbf{La apertura a 61 agrupaciones confirma que el crecimiento del PIB es aún más desigual cuando se mira con mayor detalle.} Entre las agrupaciones de mayor crecimiento aparecen Artes, entretenimiento y otros servicios (9,09\%); Transporte aéreo (6,75\%); Financieras y seguros (5,95\%); Salud y servicios sociales (5,11\%). En contraste, las mayores contracciones se observan en Minerales metalíferos (-2,32\%); Cuero, calzado y artículos de viaje (-2,28\%); Carbón (-1,98\%); Apoyo a otras actividades mineras (-0,75\%). Esta evidencia no reemplaza la medición de productividad laboral de las doce agrupaciones, pero sí permite ubicar con mayor precisión las actividades que explican el dinamismo o el rezago del PIB.
